\SourceLeaf{254}%
\Needspace{7\baselineskip}
\begin{center}
\textsc{Gauss} an \textsc{Schumacher}. G\"ottingen, 8. Februar 1846.
\end{center}

\mbox{.\ .\ .\ .\ .\ .\ .} Der Unterschied zwischen Rechts und Links l\"asst sich aber nicht definiren, sondern nur vorzeigen, so dass es damit eine \"ahnliche Bewandtniss hat, wie mit S\"uss und Bitter. Omne simile claudicat aber; das letztere gilt nur f\"ur Wesen, die Geschmacksorgane haben, das erstere aber f\"ur alle Geister, denen die materielle Welt apprehensibel ist; zwei solche Geister aber k\"onnen sich \"uber Rechts und Links nicht anders unmittelbar verst\"andigen, als indem Ein und dasselbe materielle individuelle Ding eine Br\"ucke zwischen ihnen schl\"agt, ich sage unmittelbar, da auch A sich mit Z verst\"andigen kann, indem zwischen A und B eine materielle Br\"ucke, zwischen B und C eine andere u.\,s.\,w. geschlagen werden, oder worden sein kann. Welche Geltung diese Sache in der Metaphysik hat, und dass ich darin die schlagende Widerlegung von \textsc{Kants} Einbildung finde, der Raum sei BLOSS die Form unserer \"aussern Anschauung, habe ich succinct in den G\"ottingischen Gelehrten Anzeigen 1831, S. 635 ausgesprochen. \mbox{.\ .\ .\ .\ .\ .\ .}

\par\bigskip

\Needspace{7\baselineskip}
\begin{center}
[\textsc{Gerling} \emph{an} \textsc{Gauss}. \emph{Marburg, 20. Juni 1846.}]
\end{center}

\{\mbox{.\ .\ .\ .\ .\ .\ .\ .} Zum erstenmal kam ich in den Fall, eine \textls{rechts}gewundene Schraube f\"ur den Druck [der neuen Ausgabe des »Lorenz«] definiren zu m\"ussen. Ich weiss nicht besser als so zu sagen: so gewunden, dass jeder Punkt der vertical stehenden Schraube ausser der Axe in die Tiefe nur gelangen kann, wenn er von Ost durch S\"ud nach West geht. --- Gibt es eine bessere Definition, so verpflichten Sie mich sehr durch die Mittheilung. \mbox{.\ .\ .\ .\ .\ .\ .}\}

\par\bigskip
